%!TeX root=../tese.tex
%("dica" para o editor de texto: este arquivo é parte de um documento maior)
% para saber mais: https://tex.stackexchange.com/q/78101

% As palavras-chave são obrigatórias, em português e em inglês, e devem ser
% definidas antes do resumo/abstract. Acrescente quantas forem necessárias.
\palavraschave{grafos, transporte público, ônibus, São Paulo, GTFS, mobilidade urbana, eficiência }

\keywords{graphs, public transportation, buses, São Paulo, GTFS, urban mobility, efficiency}

% O resumo é obrigatório, em português e inglês. Estes comandos também
% geram automaticamente a referência para o próprio documento, conforme
% as normas sugeridas da USP.
\resumo{
A cidade de São Paulo enfrenta diariamente desafios complexos de mobilidade urbana, com milhões de pessoas dependentes do transporte coletivo. Este trabalho tem como objetivo avaliar a eficiência da rede municipal de ônibus operada pela SPTrans e identificar potenciais pontos de melhoria em sua estrutura, por meio da modelagem matemática de redes. Utilizando dados do GTFS (General Transit Feed Specification) e séries temporais de passageiros transportados entre 2015 e 2025, incluindo o período da pandemia de Covid-19, foi realizada a integração entre informações espaciais das rotas e dados demográficos do Censo 2022 do IBGE, considerando também as subdivisões territoriais da cidade. A metodologia adotada combina técnicas de análise estatística e de redes, com a criação de grafos por meio da biblioteca NetworkX, permitindo avaliar métricas de centralidade e cobertura territorial da malha de transporte coletivo. Dentre os principais resultados, destacam-se a percepção das reestruturações do sistema após 2018, com quedas nas métricas de centralidade, indicando uma rede menos concentrada e mais distribuída, com menor dependência de poucos nós críticos, e a observação da reestruturação das rotas da Zona Oeste após inauguração da Estação Vila Sônia, em 2022, trazendo ganhos de tempo, mas também desafios de custo para os passageiros. Quanto às zonas da cidade, observamos uma maior quantidade de rotas e paradas na zona Leste e Sul, com alta demanda populacional e intenso uso das linhas de ônibus, e um planejamento mais intenso do transporte nas zonas Oeste e Centro, liderando em quantidade de rotas novas. Em relação ao perfil dos passageiros, observou-se uma redução do pagamento em dinheiro, aumento no uso de gratuidades e recuperação gradual da demanda após uma queda abrupta em 2020, ainda não atingindo os níveis pré-pandemia. Em relação à eficiência, identificamos melhorias, com avanços sociais e tecnológicos, e pontos de atenção em regiões superpopuladas e que tiveram um aumento no custo do transporte. Por fim, foi desenvolvido um painel interativo para visualização pública dos resultados, promovendo transparência e apoio ao planejamento baseado em evidências. Os resultados contribuem para uma compreensão mais aprofundada da estrutura e eficiência do transporte público paulistano, oferecendo subsídios para o planejamento urbano e a formulação de políticas de mobilidade mais equilibradas e inclusivas.


}

\abstract{
The city of São Paulo faces complex urban mobility challenges on a daily basis, with millions of people relying on public transportation. This study aims to assess the efficiency of the municipal bus network operated by SPTrans and to identify potential areas for improvement in its structure through mathematical network modeling. Using GTFS (General Transit Feed Specification) data and time series of transported passengers from 2015 to 2025, including the period of the Covid-19 pandemic, spatial information on routes was integrated with demographic data from the 2022 IBGE Census, also taking into account the city’s territorial subdivisions. The adopted methodology combines statistical and network analysis techniques, with the creation of graphs using the NetworkX library, allowing the evaluation of centrality metrics and territorial coverage of the public transport network. Among the main results, the study highlights the system’s restructurings after 2018, with declines in centrality metrics indicating a less concentrated and more distributed network, with reduced dependence on a few critical nodes, as well as the restructuring of routes in the Western Zone following the opening of Vila Sônia Station in 2022, which brought time savings but also cost challenges for passengers. Regarding the city’s regions, a greater number of routes and stops were observed in the Eastern and Southern Zones, where population density and bus usage are high, while transport planning efforts were more intense in the Western and Central Zones, which led in the creation of new routes. Concerning passenger profiles, a decrease in cash payments, an increase in the use of free passes, and a gradual recovery in demand were observed after a sharp decline in 2020, still not reaching pre-pandemic levels. In terms of efficiency, improvements were identified, with social and technological advances, alongside points of concern in overpopulated areas that experienced an increase in transportation costs. Finally, an interactive dashboard was developed for public visualization of the results, promoting transparency and supporting evidence-based planning. The findings contribute to a deeper understanding of the structure and efficiency of São Paulo’s public transportation system, providing insights for urban planning and the development of more balanced and inclusive mobility policies.

}
